% !TEX program = lualatex
\documentclass[a4paper,10pt,twocolumn]{article}
\usepackage{kaitoushu}
\renewcommand{\baselinestretch}{1.2}
\lhead{応用数学 問題集}
\problemrange{問152〜155}

\begin{document}
\thispagestyle{fancy}

\begin{center}
  {\Large \textbf{応用数学　3章　第12回　Basic}}
\end{center}
\vspace{0.5em}

\noindent\textbf{\large【Basic】}
\vspace{0.3em}

\mondai{152}{
次の関数のフーリエ変換を求めよ．
$f(x) = \begin{cases} e^{-x} & (x \geq 0) \\ 0 & (x < 0) \end{cases}$
\hfill {\scriptsize [教p.92(問1)]}
}

\mondai{153}{
次の関数のフーリエ変換を求めよ．
(1) $f(x) = \begin{cases} 2 & (-3 \leq x \leq 0) \\ 0 & (x < -3,\ x > 0) \end{cases}$
(2) $g(x) = \begin{cases} x & (0 \leq x \leq 1) \\ 0 & (x < 0,\ x > 1) \end{cases}$
\hfill {\scriptsize [教p.92(問2)]}
}

\mondai{154}{
問題152の関数にフーリエの積分定理を適用して，次の等式を証明せよ．
(1) $\dfrac{1}{2\pi}\displaystyle\int_{-\infty}^{\infty} \dfrac{1-iu}{1+u^2}\,e^{iux}\,du = \begin{cases} e^{-x} & (x>0) \\ \frac{1}{2} & (x=0) \\ 0 & (x<0) \end{cases}$
(2) $\displaystyle\int_{-\infty}^{\infty} \dfrac{1}{1+u^2}\,du = \pi$
\hfill {\scriptsize [教p.94(問3)]}
}

\mondai{155}{
次の関数のフーリエ正弦変換を求めよ．
$f(x) = \begin{cases} -x-1 & (-1 \leq x < 0) \\ -x+1 & (0 < x \leq 1) \\ 0 & (x < -1,\ x = 0,\ x > 1) \end{cases}$
\hfill {\scriptsize [教p.95(問4)]}
}

\end{document}